\title{cover letter}
%
% See http://texblog.org/2013/11/11/latexs-alternative-letter-class-newlfm/
% and http://www.ctan.org/tex-archive/macros/latex/contrib/newlfm
% for more information.
%
\documentclass[11pt,stdletter,orderfromtodate,sigleft]{newlfm}
\let\geometry\relax
\usepackage{blindtext, xfrac, animate, hyperref, pxfonts, geometry}

  \setlength{\voffset}{0in}

\newlfmP{dateskipbefore=0pt}
\newlfmP{sigsize=20pt}
\newlfmP{sigskipbefore=10pt}

\newlfmP{Headlinewd=0pt,Footlinewd=0pt}

\newcommand{\journal}{Nature Communications}
\newcommand{\articletype}{Article}
\newcommand{\myTitle}{Delayed recountings preserve but simplify the semantic geometry of earlier recountings}
\newcommand{\corresponding}{Jeremy R. Manning}


\namefrom{\vspace{-0.3in}\corresponding}
\addrfrom{
  Dartmouth College\\
  Department of Psychological and Brain Sciences\\
  HB 6207 Moore Hall\\
  Hanover, NH, 03755\\
  Jeremy.R.Manning@Dartmouth.edu}
\addrto{}
\dateset{\today}


\greetto{To the editors of \textit{\journal}:}



\closeline{Sincerely,}

\begin{document}
\begin{newlfm}

  We have enclosed our manuscript entitled \textit{\myTitle} to be considered
  for publication as an \textit{\articletype}. In our manuscript, we develop and test a mathematical framework,
  based on natural language processing models (text embeddings), for characterizing how
  people's recountings of an experienc change over time.

  We asked participants in our study to watch an episode of a television show,
  then recount it immediately, and then recount it again a week later. As in
  several of our prior studies (several of which are published in other Nature-family journals), we embedded sliding windows of detailed
  (human-generated) annotations of the episode, and of each recounting, into 
  a shared semantic space. This yields a ``trajectory'' through the embedding space
  for the episode and each recounting.

  We report two main findings. First, participants' delayed recountings resemble their own
  immediate recountings more closely than they resemble the episode they actually watched,
  and also more closely than they resemble other participants' immediate recountings. This
  suggests that what they are actually recounting after a week is best described as
  their own \textit{earlier recounting} of the episode, as opposed to either the
  actual original experience of watching the episode, or to a more general shared distillation
  of the episode that is common across participants.

  Our second finding concerns the specific changes (within person)
  between the initial (immediate) recounting and the recounting a week later. As one might expect,
  and as any reasonable memory model would predict,
  the delayed recountings capture less detail than the immediate recountings (i.e., there is some forgetting or loss of information over time).
  The interesting question is which \textit{particular} details are lost, and which are preserved.
  We applied an algorithm, originally developed for drawing cartographic maps at different scales (and with different levels of detail)
  to the trajectories of the immediate recountings. We used the algorithm to assign a score to each
  point along the trajectory that reflects how much removing that point would change the overall shape
  of the trajectory. We found that points that contributed little to the overall trajectory shapes
  were the very same points most likely to be ``dropped'' from the delayed recountings.
  In other words, delayed recountings appear to preserve the overall shape of the immediate recountings'
  paths through semantic space in a way that loses minimal information about its overall structure.

  Our work suggests that when we recount our experiences to other people, what
  we may be doing is ``leading'' them on a journey through a semantic
  representation space (approximated here using text embeddings). But just like
  navigating through the woods using hiking trail markers, not all markers are
  equally useful. It's necessary to put in a marker at ambiguous locations
  (forks in the path, overgrown parts of the trail). But markers at obvious
  locations (or, for highway navigation, unhelpful GPS ``continue straight!''
  instructions, given while already on a highway straightaway) are less useful.
  If there are "extra waypoints" in someone's recounting --- i.e., events
  inferable from their neighbors --- then we can simplify our recountings by
  dropping them, without losing much of the overall structure of the narrative.
  
We expect that our work will be of broad interest to the readership of \textit{\journal}, including
cognitive psycholoogists and cognitive neuroscientists, social psychologists, and natural language
processing researchers, among others. Thank you for considering this manuscript, and we
hope you will find it suitable for publication in \textit{\journal}.


\end{newlfm}
\end{document}
